Пусть есть два многочлена и они заданы арифметическими схемами (тоже самое что и булевы схемы, но вместо логических операций стоят арифметические). Нужно проверить эти два многочлена на равенство. Раскрытие всех скобок может занять слишком долгое время: например, произведение $n$ скобок по $2$ слагаемых после раскрытия скобок может стать экспоненциальной длины. Вместо этого используется следующая идея: посчитаем значения многочленов в случайной точке по некоторому модулю. \footnote{Без модуля возникло бы две проблемы. Во-первых, равномерного распределения на натуральных числах не бывает. Во-вторых, натуральные числа могут быть сколь угодно большими, а алгоритм должен быть полиномиальным.} Если значения разные, то многочлены тоже заведомо разные. Если же значения одинаковые, то с высокой вероятностью и многочлены одинаковые. Ниже приводится точный алгоритм и оценка вероятности его ошибки.

\begin{lstlisting}[escapeinside={(*}{*)}]
    (*Вход: Две арифметические схемы $P$ и $Q$ от $n$ переменных, задающие многочлены $p$ и $q$ степени не выше $d$*)

    (*Результат: $1$, если $p = q$ *)
            (*$0$ с вероятностью $\ge 1 - \varepsilon$, если $p \neq q$*)
            
    1   (*$N = \frac{nd}{\varepsilon} $*);
    2   (*$m$ ~--- простое число больше $N$*); // (* Можно выбрать детерминированным алгоритмом или вероятностным *)
    3   (*$x$ ~--- случайный элемент $(\Z_m)^n$*);
    4   if p(x) == q(x) mod m:
    5       return 1
    6   return 0
\end{lstlisting}

\begin{theorem} Алгоритм работает правильно, то есть возвращает заявленный результат
\end{theorem}

\begin{proof}
Если $p=q$, то для любого $x$ выполнено $p(x)=q(x)$, тем более $p(x) \equiv q(x) \pmod m$, поэтому алгоритм выдаст $1$. Если же $p \neq q$, то нужно доказать, что доля таких $x$, что $p(x) \equiv q(x) \pmod m$, не превышает $\varepsilon$. Мы докажем, что доля таких $x$, что $p(x) \not\equiv q(x) \pmod m$, не меньше $(1-\frac{d}{m})^n$, что больше $1-\varepsilon$ в силу выборов параметров и неравенства Бернулли ($(1+x)^n \ge 1 + nx$)

Доказывать будем индукцией по $n$. При $n=1$ утверждение следует из основной теоремы алгебры: у ненулевого многочлена $p-q$ степени не выше $d$ над полем $\Z_m$ не больше $d$ корней. Поэтому доля таких $x$, что $p(x) \equiv q(x) \pmod m$ не превышает $\frac{d}{m}$.

Пусть для некоторого $n$ утверждение уже доказано. Докажем для $n+1$. Пусть $x=(x_0,y) \in \Z_m \times (\Z_m)^n$. Разложим многочлены $p$ и $q$ по степеням $x_0$. Хотя бы при одной степени коэффициенты будут различными, поэтому по предположению индукции с вероятностью не меньше $(1-\frac{d}{m})^n$ они останутся разными при подстановке случайного $y$. Таким образом после подстановки $p(x_0,y)$ и $q(x_0,y)$ останутся разными многочленами от $x_0$ и с условной вероятностью не меньше $1-\frac{d}{m}$ при подстановке случайного $x_0$ дадут разные значения. Таким образом, полная вероятность будет не меньше $(1-\frac{d}{m})^{n+1}$ \end{proof}
